\chapter{项目的最小单位不是“代码”，是“实验”}
\label{ch:experiment_unit}

如果你把项目的最小单位理解成“代码”，你就会自然地用 Git 来组织一切：提交、分支、合并、回滚。
这个习惯在工程里很好，但在研究里经常失灵：\textbf{研究的产出不是某段代码，而是一个可被验证的结论。}

结论来自实验；实验由一组可追溯的输入条件生成，并产出可复查的输出。
所以本书把“实验”定义为最小单位，并围绕它组织目录结构、日志、自动化与验收标准。

\section{实验的定义：把“我改了点东西”变成可比对象}
\textbf{一个实验 = 代码版本 + 配置 + 数据版本 + 环境 + 输出 + 指标。}

你可以把它翻译成一句更落地的话：\textbf{任何一次跑出来的结果，都必须能回答“它到底是怎么来的”。}

\subsection{必须能回答的 6 个问题}
\begin{enumerate}
  \item \textbf{用的是什么代码？}（commit hash / 是否 dirty）
  \item \textbf{用的是什么配置？}（config 文件路径 + 展开后的最终参数）
  \item \textbf{用的是什么数据？}（数据版本/哈希/切分方式）
  \item \textbf{用的是什么环境？}（Python/依赖/驱动/关键硬件摘要）
  \item \textbf{产物在哪里？}（模型、日志、预测、缓存）
  \item \textbf{指标是什么、如何算的？}（评估脚本 + 后处理口径）
\end{enumerate}

\paragraph{案例：AI 生成“看似完整”的代码，但无法运行。}
跨语言迁移/“翻译代码”往往会出现 API 幻觉：函数看似合理却根本不存在，最后只能逐文件重写（见案例~\ref{case:code_translation_hallucination}）。
把“实验”当作最小单位的意义就在于：你验收的不是文件数量，而是\textbf{可运行的产物 + 可对照的指标}。

\section{实验对象模型：把研究过程拆成稳定的“名词”}
为了让流程可组合、工具可自动化，我们用一组稳定对象来描述研究过程：
\begin{itemize}
  \item \textbf{config：}一次运行的所有关键参数（可序列化、可覆盖、可 diff）。
  \item \textbf{dataset：}数据版本与切分；理想情况是可用哈希或 manifest 证明。
  \item \textbf{run：}一次具体执行（含 run\_id、时间戳、commit、config、seed、日志）。
  \item \textbf{artifact：}run 的产物（模型权重、预测结果、缓存特征、评估中间结果）。
  \item \textbf{report：}面向人类的总结（图表、表格、结论句、失败分析）。
\end{itemize}

这些对象之间的关系可以一句话说清：
\begin{quote}
config + dataset + env 经过代码执行，生成 run；run 产出 artifact；report 基于 artifact 与指标。
\end{quote}

\section{run\_id：让每一次运行都“可指代”}
在论文阶段你会频繁说：
“我们用某个设置跑出来的最好结果是……”。
如果这个设置没有一个稳定的名字，你就会陷入含糊。

建议：每次运行都生成一个\textbf{唯一 run\_id}（可读 + 可排序），例如：
\begin{verbatim}
2026-02-01_0930_baseline
2026-02-01_1130_ablation_noaug
2026-02-02_0045_sweep_lr3e-4
\end{verbatim}

并把 outputs 组织成：
\begin{verbatim}
outputs/<run_id>/
  run.json
  run.md
  metrics.json
  artifacts/
\end{verbatim}

\section{一条实用原则：把“可比性”当作第一优先级}
研究的速度往往不是被训练时间拖慢，而是被\textbf{不可比}拖慢：
\begin{itemize}
  \item baseline 和你的方法用了不同评估脚本；
  \item A 实验用了不同后处理口径；
  \item B 实验换了数据切分却没记录。
\end{itemize}

所以你应当默认：\textbf{任何结论都必须来自同一条评估流水线。}

从这一章开始，后续所有章节（仓库结构、Git、DoD、日志与 AI 工作流）都会围绕“实验可追溯 + 结果可对照”来展开。